% No preamble. Local measure factors for Notebook19. Prefix pm-.
\section{Local Haar, straightening, and exponential-gauge interactions}
\label{pm-section}

Use the exact current-step chart of Section~\ref{pc-chart-section}.
The next block side $L$ is fixed, the current spacing is one,
and constants are uniform in the preceding homogeneous depth $M$
and compatible periodic volume. Backgrounds $V=e^{ia}$ range
over a sufficiently small product ball in the independent
current link coordinates $a(b)$. A slightly larger fixed ball
is used for Cauchy estimates. This section localizes the measure
factors; the classical and previously integrated action
remainders are separate.

Write $A_{c,p}=D_{p(c)}F_V(0;c)=L^{1-d}\pi_c(V)$ and define
the square ambient straightening map
\begin{equation}\label{pm-ambient-map}
 \psi_V(B)=B-h_V\widetilde D_V(B),\qquad
 F_V(\psi_V(B))=A_VB,\qquad
 \Phi_V(z)=\psi_V(K_Vz).
\end{equation}
Its retained components equal those of $B$. The nonlinear
equation for its pivot $p(c)$ reads that pivot and no other
pivot. All maps and derivatives used below are local analytic
functions on fixed neighborhoods.

More explicitly, choose supports $Y_r$ of current cells for
every primitive term below. Include both its background and
retained fluctuation dependencies. There are fixed constants
$s_L,r_L,n_L$ such that
\begin{equation}\label{pm-support-bounds}
 |Y_r|\le s_L,\qquad \operatorname{diam}Y_r\le r_L,
 \qquad \ell(Y_r)\le s_Lr_L,\qquad
 \sup_x\#\{r:x\in Y_r\}\le n_L.
\end{equation}
Here $\ell$ is shortest-tree length; periodic multiplicities
are retained. Finite enlargement of the two-block chart
neighborhoods and the block-internal contour neighborhoods
gives these constants. The local implicit equations do not
propagate support between different coarse equations.

For nonnegative $a_0,\kappa$ define the primitive local norm
\begin{equation}\label{pm-local-norm}
 \|l\|_{a_0,\kappa,R}
 =\sup_x\sum_{r:x\in Y_r}
 e^{a_0|Y_r|+\kappa\ell(Y_r)}
 \sup_{V\in\mathcal D_b,\,\|z|_{Y_r}\|_\infty\le R}|l_r(V,z)|.
\end{equation}
The fluctuation supremum may be taken over complex coordinates,
so it also bounds the fixed real cutoff. In particular,
\begin{equation}\label{pm-incidence-bound}
 \|l\|_{a_0,\kappa,R}
 \le n_Le^{a_0s_L+\kappa s_Lr_L}\sup_r\|l_r\|_\infty.
\end{equation}
Adding primitive terms with the same support preserves this
bound by the triangle inequality.

\subsection{The exact ambient Jacobian and the delta normalization}
\label{pm-jacobian-section}

\begin{lemma}[Straightening determinant]\label{pm-jacobian-lemma}
On a fixed small complex chart,
\begin{equation}\label{pm-jacobian-identity}
 \det D\psi_V(B)=\prod_c
 \frac{\det A_{c,p}}
      {\det D_{p(c)}F_V(\psi_V(B);c)}.
\end{equation}
Each ratio is analytic, equals one at $B=0$, and has fixed
local dependence.
\end{lemma}
\begin{proof}
Order the variables as retained coordinates followed by
pivots. The retained rows of $D\psi$ are $[I,0]$.
The pivot-to-pivot block is block diagonal, because each local
implicit equation reads only its own input pivot. Differentiating
$F_V(\psi_V(B))=A_VB$ in that pivot gives
\[
 D_{p(c)}F_V(\psi_V(B);c)
 D_{B(p(c))}\psi_V(B)(p(c))=A_{c,p}.
\]
Taking the determinant of this block triangular matrix proves
\eqref{pm-jacobian-identity}. At zero, $DF_V=A_V$ and
$D\psi=I$. All local blocks stay invertible on a smaller
fixed domain.
\end{proof}

Let $\operatorname{Tr}_{\mathfrak g}$ denote the ordinary
matrix trace in color-coordinate space. The logarithm of the
normalized $c$ factor is
\begin{equation}\label{pm-jacobian-log}
 j_c^{\rm str}(V;B)=
 -\operatorname{Tr}_{\mathfrak g}\log\left[
 I+A_{c,p}^{-1}
 \{D_{p(c)}F_V(\psi_V(B);c)-A_{c,p}\}\right],
\end{equation}
with branch zero at $B=0$. This is a determinant identity:
the diagonal Jacobian block itself is $(DF_p)^{-1}A_p$,
and no commutativity of color matrices is asserted.
Local analyticity and the logarithm series imply, for a fixed
$\rho_L>0$ and $|g|R<\rho_L$,
\begin{equation}\label{pm-jacobian-bound}
 \|j^{\rm str}(V;K_Vgz)\|_{a_0,\kappa,R}
 \le C_{L,a_0,\kappa}|g|R.
\end{equation}
Indeed, the matrix perturbation inside the logarithm is at
most $C_L|g|R$, which may be made at most $1/2$.

The determinant in Lemma~\ref{pm-jacobian-lemma} belongs to
the square ambient change of variables, not to the rectangular
derivative $D\Phi_V$. To identify its measure role, for real
backgrounds and test functions supported within the chart,
\begin{equation}\label{pm-delta-reduction}
 \begin{split}
 \int dX\,\delta(F_V(X))f(X)
 &=\mathcal D(V)\int dz\,
 \exp\left\{\sum_cj_c^{\rm str}(V;K_Vz)\right\}
 f(\Phi_V(z)),\\
 \mathcal D(V)&=\prod_c(\det A_{c,p})^{-1}.
 \end{split}
\end{equation}
First set $X=\psi_V(B)$ and then integrate the linear
constraint $A_VB=0$ over pivots. The pivot determinants have
the positive orientation of $L^{1-d}I$ on this real
neighborhood. A group delta at identity contributes only its
fixed exponential-coordinate normalization at $F=0$.
After retained variables are scaled by positive real $g$,
the additional factor $g^{\dim z}$ is field independent.
Complex-background formulas are holomorphic continuations
of the local densities; no complex distributional delta is used.

The factor $\mathcal D(V)$ is independent of fluctuations
but can depend on the prescribed next-scale field through $V$.
It is kept outside the fluctuation interaction, not discarded.
Its normalized logarithm is local:
\begin{equation}\label{pm-background-determinant}
 \log\frac{\mathcal D(V)}{\mathcal D(1)}
 =-\sum_c\operatorname{Tr}_{\mathfrak g}
       \log(A_{c,p}A_{1,c,p}^{-1}).
\end{equation}
Each term is $O(\|a\|_\infty)$ and has no small coupling
factor $g$.

\subsection{The relative Haar density}
\label{pm-haar-section}

Let $j(X)$ be the pullback density of Haar measure under
$X\mapsto e^{iX}$, divided by its value at zero. Near zero
it is nonvanishing and real analytic, with $j(0)=1$.
Haar invariance under inversion gives $j(-X)=j(X)$ on
the real form and hence by analytic continuation. Its local
logarithm therefore satisfies
\begin{equation}\label{pm-haar-even}
 \log j(0)=D\log j(0)=0,\qquad |\log j(X)|\le C|X|^2.
\end{equation}
This uses the positive real density and then its analytic
branch; it does not complexify an absolute determinant.
Translation by $V$ preserves Haar measure, so the relative
parametrization $e^{iX}V$ has the same density.

The exact Haar logarithms after scaling are
\begin{equation}\label{pm-haar-log}
 j_b^{\rm H}(V;gz)=\log j(\Phi_V(gz)(b)).
\end{equation}
On retained bonds the argument is exactly $gz(b)$; on pivots
it is the local chart function. From
\eqref{pm-haar-even}, the uniform local chart bound, and
\eqref{pm-incidence-bound},
\begin{equation}\label{pm-haar-bound}
 \|j^{\rm H}(V;gz)\|_{a_0,\kappa,R}
 \le C_{L,a_0,\kappa}|g|^2R^2.
\end{equation}
Any unnormalized value $j(0)$ is field independent and remains
in the overall normalization. A generally nonzero quadratic
Haar term remains in the interaction with coefficient $g^2$.

\subsection{The exact retained-coordinate gauge remainder}
\label{pm-gauge-section}

For a root-to-site pair $r=(y,x)$ define
\begin{equation}\label{pm-gauge-primitive}
 f_r(V;z)=1-\tfrac12\tau\left[
 P_{e^{i\Phi_V(z)}V}(y,x)
 +P_{e^{i\Phi_V(z)}V}(y,x)^{-1}\right].
\end{equation}
The source contours are internal to blocks and contain no
pivots. By \eqref{pc-gauge-pivot-independence}, this is
exactly the same expression with $\Phi_V(z)$ replaced by
$E_fz$, where $E_f$ inserts zero at pivots. Consequently
it is a local function of unchanged retained coordinates.

On the full small product background domain define
\begin{equation}\label{pm-gauge-remainder}
 \begin{split}
 v_r^{\rm G}(V;g,z)=-g^{-2}\bigl[
 &f_r(V;gz)-f_r(V;0)-gD_zf_r(V;0)z\\
 &-\tfrac12g^2D_z^2f_r(V;0)[z,z]\bigr].
 \end{split}
\end{equation}
The singularity at $g=0$ is removable, with value zero.
The finite-dimensional Taylor remainder gives
\begin{equation}\label{pm-gauge-bound}
 \|v^{\rm G}(V;g,z)\|_{a_0,\kappa,R}
 \le C_{L,a_0,\kappa}|g|R^3,
 \qquad |g|R<\rho_L.
\end{equation}
At a physical background $P_V=1$, both $f_r(V;0)$ and
$D_zf_r(V;0)$ vanish. Thus this is precisely the logarithmic
gauge interaction after its quadratic contribution has been
included in the precision. Elsewhere it is the exact
third-and-higher remainder after constant and linear terms
are also separated. The full-domain definition supplies
background holomorphy without treating $P_V=1$ as a product
domain.

\subsection{Local coefficient and multiplier bounds}
\label{pm-coefficients-section}

\begin{theorem}[Local measure interaction]\label{pm-measure-theorem}
For the fixed geometry and product background domain above,
the logarithmic interactions
$j^{\rm str}(V;K_Vgz)$, $j^{\rm H}(V;gz)$ and
$v^{\rm G}(V;g,z)$ are holomorphic and have primitive
supports satisfying \eqref{pm-support-bounds}. For
$|g|R<\rho_L$, their sum, regrouped by supports $Y$, obeys
\begin{equation}\label{pm-combined-bound}
 \sup_x\sum_{Y\ni x}e^{a_0|Y|+\kappa\ell(Y)}
 \sup|v_Y(V;g,z)|
 \le C_{L,a_0,\kappa}
   (|g|R+|g|^2R^2+|g|R^3).
\end{equation}
The constants are independent of $M$ and periodic volume.
If the right side is at most one, the same bound multiplied
by $e$ holds with $|v_Y|$ replaced by $|e^{v_Y}-1|$.
For $R\ge1$ and $|g|R\le1$, it is at most
$C_{L,a_0,\kappa}|g|R^3$.
\end{theorem}
\begin{proof}
Combine \eqref{pm-jacobian-bound},
\eqref{pm-haar-bound} and \eqref{pm-gauge-bound} and group
equal supports. The local coefficient assertion can also
be seen directly from the convergent expansions
\[
 \begin{split}
 j_c^{\rm str}(V;K_Vgz)&=\sum_{n\ge1}g^n s_{c,n}(V;z),\\
 j_b^{\rm H}(V;gz)&=\sum_{n\ge2}g^n h_{b,n}(V;z),\\
 v_r^{\rm G}(V;g,z)&=\sum_{n\ge3}g^{n-2}v_{r,n}(V;z).
 \end{split}
\]
Each coefficient is homogeneous of degree $n$ in $z$,
holomorphic in only the background coordinates of its support,
and has local norm at most $C_0(C_1R)^n$ by Cauchy on a
fixed fluctuation disc. Background Cauchy estimates on the
smaller product domain preserve the same supports. These are
anchored coefficient bounds, rather than an extensive bound
on the full sum. Finally $|e^v-1|\le e|v|$ for $|v|\le1$.
\end{proof}

The support-cardinality factor $K^{|Y|}$ in the overlapping
interaction theorem is absorbed by increasing $a_0$ by
$\log\max(1,K)$. This is a fixed-$L$ statement, not a
claim of uniform constants when $L$ or the cell dimension is
allowed to increase. For a logarithmically growing cutoff
$R(g)$, these measure interactions tend to zero as soon as
$|g|R(g)^3\to0$.

The fixed real cutoff is the original retained window
$\|z(b)\|<R$. It remains unchanged because
$\Phi_V(gz)(b)=gz(b)$ on retained bonds. Additional pivot,
contour or reconstructed-field indicators retain their exact
pullbacks and are not covered by this theorem. CMP116,
p.~12 immediately before (2.3), explicitly uses the
pivot-removed bond set; its (2.14), p.~15, keeps this retained
cutoff fixed during Gaussian deformation. CMP109,
pp.~266--268, (2.5) and (2.9)--(2.13), supplies the gauge
and pivot setup. The unreadable product subscript of
CMP109 (2.9) in the accessible OCR is not presented as a
visually verified formula.

These factors are strictly local in independent current
backgrounds $V$. Substituting $V_*(W)$ does not give them
strict finite support in $W$. The full classical remainder,
the previously integrated interaction, the external factor
\eqref{pm-delta-reduction}, and any extra cutoff indicators
remain distinct obligations.
